\subsection{Existing Optimisation Approaches}

Current approaches to optimising UAV-LoRaWAN networks generally fall into two categories:
static mathematical optimisation and energy-focused path planning.

In the domain of mathematical optimisation, Zorbas et al.~\cite{zorbas2025revisiting} propose
a Mixed Integer Linear Programming (MILP) framework to determine the optimal SF distribution
that maximises packet reception probability. While this method provides a theoretically optimal
allocation for static networks, it has two critical limitations for the mobile-gateway setting:
first, MILP scales exponentially with the number of binary SF assignment variables, becoming
intractable for $N > 20$ sensors; second, and more fundamentally, it assumes channel conditions
are known and stationary at solve time, an assumption that fails continuously as the UAV moves
and the EMA-ADR control loop propagates outdated RSSI estimates across the network. The authors
themselves acknowledge this limitation, concluding that reinforcement learning is the appropriate
tool for the dynamic case~\cite{zorbas2025revisiting}.

In the domain of path planning, Xiong et al.~\cite{xiong2023flyinglora} introduce
``FlyingLoRa,'' a scheme designed to minimise energy consumption by jointly optimising the
UAV's 3D trajectory and transmission scheduling. Although this represents the state-of-the-art
in energy-efficient LoRa data collection, FlyingLoRa has three notable limitations: it
optimises trajectory offline using iterative algorithms that require full knowledge of sensor
locations and link budgets before the mission begins, it omits ADR dynamics entirely (treating
each sensor's data rate as fixed), and it does not model the Capture Effect---assuming mutual
packet destruction on any intra-SF collision. As a result, FlyingLoRa cannot exploit the
positioning-induced SF diversity that the protocol-aware DRL agent learns to generate.

A parallel line of research addresses UAV trajectory optimisation as an energy-efficiency
problem independent of the application-layer protocol. Zeng and
Zhang~\cite{zeng2017energy} derive a theoretical propulsion energy model for fixed-wing
UAVs and show that jointly optimising the flight trajectory and communication scheduling
yields significant energy efficiency gains; this framework motivates the energy-consumption
term in our reward function. However, the Zeng--Zhang model is derived for fixed-wing
platforms under the assumption of LoS-only links, making it inapplicable to rotary-wing
UAVs (where translational lift inverts the move/hover power relationship) and environments
with ground reflection. More broadly, Zeng et al.~\cite{zeng2016wireless} and
Mozaffari et al.~\cite{mozaffari2019tutorial} survey the channel characteristics and
deployment challenges unique to UAV-aided wireless systems, establishing the channel
geometry assumptions---in particular the dominance of LoS links at low altitude---on which
the Two-Ray Ground Reflection model in this work is predicated. A key gap across all of
these works is that none jointly models the physical channel \emph{and} the MAC-layer ADR
feedback loop, meaning their conclusions about optimal positioning do not account for
the SF convergence dynamics that dominate throughput in dense deployments.

At the intersection of deep reinforcement learning and UAV communications, Challita
et al.~\cite{challita2018interference} propose a deep RL framework for interference-aware
path planning of cellular-connected UAVs, demonstrating that DRL agents can learn
collision-avoidance behaviour that static planners cannot replicate---directly analogous to
the SF-monoculture avoidance objective here. However, Challita et al.\ operate in the
cellular domain where the interference model is well-characterised and stationary; they
do not encounter the non-stationary MAC-layer dynamics (EMA-ADR, duty cycles, Capture
Effect) that make the LoRaWAN setting substantially harder.
Nemer et al.~\cite{nemer2022energyefficient} apply a distributed actor--critic approach to
maximise coverage fairness for multi-UAV networks, confirming that Jain's Fairness Index
is an appropriate target metric for DRL-based UAV controllers. A limitation of Nemer et al.\
is that their evaluation uses fixed sensor positions and a simplified channel model,
leaving open whether the fairness gains transfer to the stochastic, topology-varying
conditions that domain randomisation is designed to address.

\subsubsection{Algorithm Selection: DQN versus Policy-Gradient Methods}

Three candidate families are considered: DQN~\cite{mnih2015humanlevel}, PPO~\cite{schulman2017proximal}, and SAC~\cite{haarnoja2018soft}.
The UAV's five-action discrete space directly favours DQN: it was designed for discrete
actions and achieves sample efficiency via experience replay and a target network.
PPO handles discrete actions but is on-policy, discarding all experience after each
update --- a significant cost when episodes span 2100 steps and interactions are
computationally expensive. SAC's discrete variant~\cite{christodoulou2019soft} loses
the entropy-maximisation guarantee that motivates its exploration, offering no clear
advantage here.

Experience replay also addresses the EMA-ADR non-stationarity: by re-using transitions
from earlier training states, it provides a diverse distribution that stabilises the
Q-function as the reward landscape shifts with ADR convergence --- a property on-policy
PPO cannot replicate. Challita et al.~\cite{challita2018interference} and Luong
et al.~\cite{luong2019drl} both support DQN for discrete-waypoint UAV navigation,
motivating it as the primary algorithm with actor--critic extensions as future work
(Section~\ref{sec:future_work}).
